\chapter{Estado del arte}
\label{cap:capitulo2}

Este capítulo detalla algunos de los sistemas y soluciones más significativos en el campo de la dispensación automática de medicamentos, además de distintos métodos tecnológicos vinculados a la incorporación de sistemas \textit{IoT}, arquitecturas inteligentes y aplicaciones móviles que buscan optimizar la adherencia terapéutica.\\

Se han revisado propuestas académicas publicadas en revistas de alto impacto y también implementaciones prácticas de código abierto, para estudiar sus propiedades técnicas, beneficios, restricciones y su uso en la creación de un dispensador automático que cuente con pantalla táctil y aplicación móvil.\\

\section{A Systematic Review on Pill and Medication Dispensers from a Human-Centered Perspective}
\label{sec:systematic_review}

En el artículo presentado por \cite{Gargioni2024} se realizó una revisión sistemática sobre dispensadores automáticos de medicamentos siguiendo la metodología \textit{PRISMA (Preferred Reporting Items for Systematic Reviews and Meta-Analyses}). El estudio analizó soluciones publicadas en bases de datos como \textit{Scopus} y \textit{PubMed} con el objetivo de responder a la siguiente cuestión: qué tipos de soluciones automatizadas se han propuesto en la literatura científica para mejorar la adherencia terapéutica.\\

Después de un proceso de selección, como se muestra en la Figura \ref{fig:prisma_diagram}, se redujeron los 426 artículos iniciales a 56 trabajos finales. La clasificación de las soluciones estudiadas se realizó según su estado de desarrollo (ya sea prototipo, estudio conceptual o producto comercial) y se examinó desde cuatro perspectivas: la tecnología empleada, la arquitectura hardware/software, el grado de participación de los distintos actores (tales como cuidadores, pacientes y personal sanitario) y los tipos de interacción.\\

\begin{figure} [H]
  \begin{center}
    \includegraphics[width=\textwidth]{figs/prisma_diagram.png}
  \end{center}
  \caption{Diagrama PRISMA del proceso de selección de artículos. Adaptado de \cite{Gargioni2024}.}
  \label{fig:prisma_diagram}
\end{figure}

Un aspecto a destacar fue su perspectiva centrada en el usuario. A diferencia de otras revisiones que solo se centraban en los aspectos técnicos, este estudio también analizó variables como la aceptabilidad, la usabilidad y el contexto de uso, especialmente en pacientes con enfermedades crónicas y personas mayores.\\

El análisis reveló que la mayoría de los sistemas incorporaban funciones como recordatorios por medio de señales acústicas o luminosas, conexión con aplicaciones para móviles y conectividad a la nube para supervisión desde lejos. No obstante, también se identificaron deficiencias relevantes en la integración con los sistemas de salud oficiales, en el manejo de terapias complejas y en la adaptación a personas con deterioro cognitivo.\\

\

Las ventajas de este estudio fueron las siguientes:

\begin{itemize}
    \item Aplicación rigurosa de metodología PRISMA.
    \item Clasificación clara de las soluciones según su estado de desarrollo.
    \item Análisis conjunto de arquitectura técnica e interacción usuario-sistema.
    \item Enfoque centrado en la aceptabilidad y usabilidad.\\
\end{itemize}

Las desventajas identificadas en las soluciones revisadas fueron las siguientes:

\begin{itemize}
    \item Falta de estandarización en la medición de la adherencia terapéutica.
    \item Escasa integración con sistemas sanitarios nacionales.
    \item Coste elevado en varias propuestas comerciales.
    \item Limitada consideración de la experiencia de usuario en población mayor.
\end{itemize}

\

\section{mPillBox: An Open Framework for a Smarter Pillbox}
\label{sec:mPillBox}

En el trabajo presentado por \cite{alduaij2024} se introduce \textit{mPillBox}, un \textit{framework} abierto diseñado para transformar un pastillero convencional en un sistema inteligente mediante la incorporación de tecnologías \textit{IoT}. Esta propuesta, en comparación con otros sistemas que sugerían una reestructuración total del aparato físico, se fundamentó en una arquitectura modular, la cual posibilitó incorporar funcionalidades de avanzada a estructuras ya existentes.\\

El sistema se puso en marcha sobre una \textit{Raspberry Pi 4 Model B}, la cual actuó como unidad de coordinación y procesamiento central. La arquitectura incorporó módulos de comunicación para enlazarse con dispositivos externos y aplicaciones móviles, además de sensores ambientales para monitorizar situaciones como la humedad y la temperatura. Además incluyó funciones como la programación de horarios para la medicación, la verificación de fechas de vencimiento y el chequeo automático de posibles interacciones entre medicamentos.\\

Como se puede observar en la Figura \ref{fig:mpillbox}, el prototipo integró físicamente la unidad de procesamiento, el módulo de visualización LCD y los elementos de conexión necesarios para la comunicación con los compartimentos del pastillero, conformando una arquitectura modular basada en componentes de bajo coste.\\

\begin{figure} [H]
  \begin{center}
    \includegraphics[width=\textwidth]{figs/mPillBox.png}
  \end{center}
  \caption{Prototipo del sistema mPillBox basado en Raspberry Pi y módulos \textit{IoT}. Adaptado de \cite{alduaij2024}.}
  \label{fig:mpillbox}
\end{figure}

Una de las características distintivas del sistema fue su perspectiva amplia y abierta. \textit{mPillBox} se diseñó como una plataforma adaptable que simplificó la incorporación de nuevos módulos y servicios, lo cual posibilitó su ajuste a diversas necesidades y contextos. Asimismo, no estaba directamente vinculado con instalaciones hospitalarias ni con servicios de suscripción, lo que disminuyó los costos y mejoró su accesibilidad.\\

El sistema también contempló el envío de alertas mediante dispositivos \textit{IoT} conectados, como notificaciones visuales o sonoras en el entorno doméstico, así como la sincronización con una aplicación móvil para la supervisión remota del tratamiento.\\

\

\

\

\

\

Las ventajas de este sistema fueron las siguientes:

\begin{itemize}
    \item Arquitectura \textit{open-source} y modular.
    \item Bajo coste en comparación con soluciones comerciales cerradas.
    \item Flexibilidad para integrar sensores adicionales.
    \item Incorporación de verificación de interacciones entre medicamentos.
    \item Independencia de infraestructuras hospitalarias externas.\\
\end{itemize}

Las desventajas de este sistema fueron las siguientes:

\begin{itemize}
    \item Requería conocimientos técnicos para su instalación y mantenimiento.
    \item No incorporó un mecanismo avanzado de dispensación física automática.
    \item Validación limitada a entorno experimental.
    \item Dependencia de conectividad para algunas funcionalidades \textit{IoT}.
\end{itemize}

\

\section{Intelligent and Assisted Medicine Dispensing Machine}
\label{sec:intelligent_and_assisted}

En el trabajo presentado por \cite{intelligent2023} se creó y diseñó una máquina dispensadora de medicamentos inteligente, enfocada en pacientes con problemas para manejar tratamientos complejos y personas mayores. La finalidad del sistema consistió en optimizar la adherencia terapéutica a través de la automatización de la dispensación y de la inclusión de mecanismos extra de seguridad.\\

La arquitectura del dispositivo se basó en un microcontrolador encargado de gestionar los compartimentos de medicación y controlar el mecanismo de liberación mediante actuadores eléctricos. El sistema permitió programar horarios específicos para cada toma, activando la dispensación de forma automática en los intervalos establecidos.\\

\begin{figure} [H]
  \begin{center}
    \includegraphics[width=\textwidth]{figs/mdm.png}
  \end{center}
  \caption{Arquitectura propuesta del sistema. Adaptado de \cite{intelligent2023}.}
  \label{fig:mdm_architecture}
\end{figure}

Como se muestra en la Figura \ref{fig:mdm_architecture}, el sistema integra los distintos módulos de control, sensores y elementos de comunicación que permiten coordinar la autenticación del usuario, la apertura de los compartimentos y la monitorización remota mediante el sistema central.\\

Un sistema de autenticación biométrica a través de un sensor de huella dactilar fue uno de los componentes distintivos del proyecto. Al permitir la validación de la identidad del usuario antes de dar el visto bueno a la dispensación, este mecanismo disminuyó el peligro de que se cometieran errores en la administración del fármaco o de que se accediera indebidamente.\\

Además, el sistema incorporó conectividad IoT para el envío de notificaciones y alertas a familiares o cuidadores en caso de omisión de la toma. La interfaz de usuario se diseñó con una pantalla LCD y botones físicos, buscando simplicidad y facilidad de uso para población de edad avanzada. Por último, el prototipo fue validado experimentalmente en un entorno controlado, demostrando la viabilidad técnica de la arquitectura propuesta.\\

Las ventajas de este sistema son las siguientes:

\begin{itemize}
    \item Integración de autenticación biométrica mediante huella dactilar.
    \item Reducción del riesgo de uso indebido del dispositivo.
    \item Automatización de la dispensación en horarios programados.
    \item Incorporación de conectividad IoT para monitorización remota.
    \item Enfoque orientado a población mayor.
\end{itemize}

Las desventajas de este sistema son las siguientes:

\begin{itemize}
    \item Incremento del coste debido a la incorporación de sensores biométricos.
    \item Mayor complejidad técnica respecto a soluciones básicas.
    \item Dependencia de conectividad para funcionalidades remotas.
    \item Validación limitada a entorno experimental.
\end{itemize}
 
\section{Metodología}
\label{sec:metodologia}

Qué paradigma de desarrollo software has seguido para alcanzar tus objetivos.

\section{Plan de trabajo}
\label{sec:plantrabajo}

Qué agenda has seguido. Si has ido manteniendo reuniones semanales, cumplimentando objetivos parciales, si has ido afinando poco a poco un producto final completo, etc.
